\problem{3.31}{Orthonormal Legendre polynomials from Gram–Schmidt}
On $[-1,1]$ take the inner product
$\braket fg=\int_{-1}^1 f(x)g(x)\dd x$.
Start with $u_0=1$, $\|u_0\|=\sqrt2$.
Since $\int x\dd x=0$, $u_1=x$,
$\|u_1\|=\sqrt{2/3}$.
For $x^2$, subtract the projection onto $1$:
\begin{equation}
u_2=x^2-\frac{\int_{-1}^1x^2\dd x}
{\int_{-1}^1 1\dd x}=x^2-\frac13.
\end{equation}
It is automatically orthogonal to $x$ by parity,
and $\|u_2\|^2=8/45$.
For $x^3$, subtract the projection onto $x$:
\begin{equation}
u_3=x^3-\frac{\int x^4\dd x}
{\int x^2\dd x}\,x=x^3-\frac35x.
\end{equation}
The other projections vanish by parity or
orthogonality, and $\|u_3\|^2=8/175$.
The orthonormal basis is therefore
\begin{align}
\boxed{e_0(x)}&=\frac1{\sqrt2},&
\boxed{e_1(x)}&=\sqrt{\frac32}\,x,\\
\boxed{e_2(x)}&=\sqrt{\frac52}\,
\frac{3x^2-1}{2},&
\boxed{e_3(x)}&=\sqrt{\frac72}\,
\frac{5x^3-3x}{2}.
\end{align}
These are $\sqrt{(2n+1)/2}\,P_n(x)$
for Legendre polynomials normalized by $P_n(1)=1$.

\problem{3.32}{Anti-Hermitian operators and their commutators}
Let $Q^\dagger=-Q$.
Then
$\avg Q^*=\bra\psi Q^\dagger\ket\psi=-\avg Q$,
so $\avg Q$ is purely imaginary (or zero).
For $Q\ket u=\lambda\ket u$ with $\|u\|=1$,
$\lambda=\avg Q$ is likewise imaginary.
For eigenvectors $u,v$ with eigenvalues
$\lambda,\mu$, respectively,
\begin{equation}
\braket{Qu}v=\lambda^*\braket uv
=-\braket u{Qv}=-\mu\braket uv.
\end{equation}
Hence $(\lambda^*+\mu)\braket uv=0$.
Since $\lambda^*=-\lambda$,
distinct eigenvalues imply
$\boxed{\braket uv=0}$.
For Hermitian $A,B$,
\begin{equation}
[A,B]^\dagger=(AB-BA)^\dagger=BA-AB=-[A,B].
\end{equation}
The same conclusion holds if both operators
are anti-Hermitian: their commutator is
again anti-Hermitian.
Finally every operator $Q$ decomposes as
\begin{equation}
\boxed{Q=A+B,\quad A=\frac{Q+Q^\dagger}{2},
\quad B=\frac{Q-Q^\dagger}{2}},
\end{equation}
with $A^\dagger=A$ and $B^\dagger=-B$.

\problem{3.33}{Three sequential projective measurements}
After measuring $A$ and obtaining $a_1$,
the state is $\psi_1=(3\phi_1+4\phi_2)/5$.
A subsequent measurement of $B$ yields
$b_1$ with probability $9/25$ or $b_2$
with probability $16/25$.
Conditional on $b_1$, the state becomes
$\phi_1$ and the next probability of
$a_1$ is $|\braket{\psi_1}{\phi_1}|^2=9/25$;
conditional on $b_2$, that probability is
$16/25$.
If the $B$ outcome is not disclosed,
the law of total probability gives
\begin{equation}
\boxed{P(a_1\text{ at the third measurement})=
\frac9{25}\frac9{25}+
\frac{16}{25}\frac{16}{25}
=\frac{337}{625}}.
\end{equation}
The result differs from the probability one
without an intervening nonselective $B$ measurement,
because that measurement destroys coherence
in the $B$ basis.

\problem{3.34}{Infinite-well energy eigenstates in momentum space}
For $0<x<a$ write
$\psi_n(x)=\sqrt{2/a}\sin(k_nx)$,
$k_n=n\pi/a$, and zero otherwise.
The Fourier transform integral is
\begin{align}
\Phi_n(p,t)&=\frac{\ee^{-\ii E_nt/\hbar}}
{\sqrt{2\pi\hbar}}\sqrt{\frac2a}
\int_0^a\sin(k_nx)\ee^{-\ii px/\hbar}\dd x\\
&=\boxed{
\frac{\ee^{-\ii E_nt/\hbar}}{\sqrt{\pi a\hbar}}
\frac{k_n[1-(-1)^n\ee^{-\ii pa/\hbar}]}
{k_n^2-(p/\hbar)^2}}.
\end{align}
Apparent poles at $p=\pm\hbar k_n$ are
removable: the numerator vanishes there too.
Writing $q=p/\hbar$, the probability density is
\begin{equation}
\boxed{|\Phi_n(p,t)|^2=
\frac{2k_n^2[1-(-1)^n\cos(qa)]}
{\pi a\hbar(k_n^2-q^2)^2}}.
\end{equation}
It is symmetric in $p$, hence $\avg p=0$.
For large $n$, the dominant momentum peaks
are near $p=\pm n\pi\hbar/a$, reflecting the
two traveling-wave constituents of the standing
energy eigenstate. Parseval's identity for
Fourier transforms gives
\begin{equation}
\avg{p^2}=
\hbar^2\int_{\mathbb R}|\psi_n'(x)|^2\dd x
=\boxed{\left(\frac{n\pi\hbar}{a}\right)^2},
\end{equation}
in agreement with Problem 2.4.
\qnote{The PDF's numerical plots use a
dimensionless momentum $pa/\hbar$ and
$n=1,2,5,10$.}

\problem{3.35}{Finite wave train and the cost of sharp edges}
Set $L=n\lambda$ with positive integer $n$,
$k_0=2\pi/\lambda$, and
$\Psi(x,0)=\ee^{\ii k_0x}/\sqrt{2L}$
on $-L<x<L$, zero elsewhere.
Its Fourier transform is
\begin{align}
\Phi(p,0)
&=\frac1{\sqrt{4\pi L\hbar}}
\int_{-L}^L\ee^{-\ii(p/\hbar-k_0)x}\dd x\\
&=\boxed{\frac1{\sqrt{\pi L\hbar}}
\frac{\sin[(p/\hbar-k_0)L]}{p/\hbar-k_0}}.
\end{align}
The position density is a rectangle of height
$1/(2L)$ and width $w_x=2L$.
The momentum density is a squared-sinc distribution
centered at $p_0=\hbar k_0$.
Its first zeros are at $p=p_0\pm\pi\hbar/L$,
so the main-lobe width is
\begin{equation}
\boxed{w_p=\frac{2\pi\hbar}{L}},\qquad
\boxed{w_xw_p=4\pi\hbar}.
\end{equation}
As $n\to\infty$, $w_x\to\infty$ and
$w_p\to0$. These widths are only
rough spread estimates, not standard deviations.
Indeed $|\Phi(p)|^2$ decays as $p^{-2}$
with nonvanishing oscillations. Therefore
$\int p^2|\Phi(p)|^2\dd p$ diverges and
$\sigma_p=\infty$. The sharp cutoff
creates jumps in $\Psi$ and invalidates a
finite kinetic-energy variance.
\qnote{One must not claim a finite Heisenberg
standard-deviation product based on the
sinc peak widths.}

\problem{3.36}{Lorentzian-amplitude wavefunction and its Fourier dual}
Take $\Psi(x,0)=A/(x^2+a^2)$ with $a>0$.
Use
\begin{equation}
\int_{-\infty}^\infty\frac{\dd x}{(x^2+a^2)^2}
=\frac{\pi}{2a^3},
\quad
\int_{-\infty}^\infty\frac{x^2\dd x}{(x^2+a^2)^2}
=\frac{\pi}{2a}.
\end{equation}
Normalization gives
$\boxed{A=\sqrt{2a^3/\pi}}$.
Parity yields $\avg x=0$, and the second
integral gives $\avg{x^2}=a^2$,
so $\boxed{\sigma_x=a}$.
The standard Fourier integral
$\int\ee^{-\ii kx}/(x^2+a^2)\dd x
=(\pi/a)\ee^{-a|k|}$ gives
\begin{equation}
\boxed{\Phi(p,0)=
\sqrt{\frac a\hbar}\ee^{-a|p|/\hbar}}.
\end{equation}
Check normalization:
$\int|\Phi|^2\dd p=
(2a/\hbar)\int_0^\infty\ee^{-2ap/\hbar}\dd p=1$.
Again $\avg p=0$, and
\begin{equation}
\avg{p^2}=\frac{2a}{\hbar}
\int_0^\infty p^2\ee^{-2ap/\hbar}\dd p
=\frac{\hbar^2}{2a^2}.
\end{equation}
Thus
\begin{equation}
\boxed{\sigma_p=\frac{\hbar}{\sqrt2a}},
\qquad
\boxed{\sigma_x\sigma_p=\frac{\hbar}{\sqrt2}
>\frac{\hbar}{2}}.
\end{equation}

\problem{3.37}{General virial theorem and the harmonic oscillator}
For $H=p^2/(2m)+V(x)$, apply
Ehrenfest's identity to the operator $xp$
(on an appropriate domain):
\begin{equation}
\frac{\dd}{\dd t}\avg{xp}=
\frac{\ii}{\hbar}\avg{[H,xp]}.
\end{equation}
Use $[H,xp]=[H,x]p+x[H,p]$
and the commutators
$[H,x]=-\ii\hbar p/m$,
$[H,p]=\ii\hbar V'(x)$:
\begin{equation}
\boxed{\frac{\dd}{\dd t}\avg{xp}
=2\avg T-\avg{xV'(x)}}.
\end{equation}
For a stationary energy eigenstate, the
expectation of a time-independent observable
is constant, hence the left side vanishes:
\begin{equation}
\boxed{2\avg T=\avg{xV'(x)}}.
\end{equation}
For $V=m\omega^2x^2/2$,
$xV'=m\omega^2x^2=2V$,
so $\boxed{\avg T=\avg V=E_n/2}$,
consistent with oscillator energy moments.

\problem{3.38}{Time to become orthogonal for two energy levels}
At $t=0$ let
$\ket{\Psi(0)}=(\ket1+\ket2)/\sqrt2$
with orthonormal energy states and
$E_1\ne E_2$.
Time evolution gives
\begin{equation}
\ket{\Psi(t)}=\frac1{\sqrt2}
(\ee^{-\ii E_1t/\hbar}\ket1+
\ee^{-\ii E_2t/\hbar}\ket2).
\end{equation}
Its overlap with the initial state is
\begin{align}
\braket{\Psi(0)}{\Psi(t)}
&=\frac12(\ee^{-\ii E_1t/\hbar}
+\ee^{-\ii E_2t/\hbar})\\
&=\ee^{-\ii(E_1+E_2)t/(2\hbar)}
\cos\left(\frac{(E_2-E_1)t}{2\hbar}\right).
\end{align}
The first orthogonality time is
$\boxed{\tau=\pi\hbar/|E_2-E_1|}$.
The energy standard deviation is
$\boxed{\sigma_H=|E_2-E_1|/2}$.
With the specified $\Delta t=\tau/\pi$,
\begin{equation}
\boxed{\sigma_H\Delta t=
\frac{|E_2-E_1|}{2}\,
\frac{\hbar}{|E_2-E_1|}
=\frac{\hbar}{2}}.
\end{equation}
The case of equal energy does not
orthogonalize and has $\sigma_H=0$.

\problem{3.39}{Oscillator matrices for position, momentum and Hamiltonian}
Let $a\ket n=\sqrt n\ket{n-1}$,
$a^\dagger\ket n=\sqrt{n+1}\ket{n+1}$.
With $\ell_x=\sqrt{\hbar/(2m\omega)}$,
$\ell_p=\sqrt{m\hbar\omega/2}$,
\begin{align}
x&=\ell_x(a+a^\dagger),&
p&=\ii\ell_p(a^\dagger-a).
\end{align}
Therefore, for nonnegative indices $j,k$,
\begin{align}
\boxed{X_{jk}}&=
\ell_x(\sqrt{k}\,\delta_{j,k-1}
+\sqrt{k+1}\,\delta_{j,k+1}),\\
\boxed{P_{jk}}&=
\ii\ell_p(\sqrt{k+1}\,\delta_{j,k+1}
-\sqrt{k}\,\delta_{j,k-1}).
\end{align}
Thus $X$ is real symmetric and tridiagonal,
and $P$ is imaginary Hermitian and tridiagonal.
Squaring in the \emph{complete infinite basis},
\begin{equation}
\frac{P^2}{2m}+\frac12m\omega^2X^2
=\frac{\hbar\omega}{4}
[(a^\dagger-a)^2(-1)+(a+a^\dagger)^2]
=\frac{\hbar\omega}{2}(aa^\dagger+a^\dagger a).
\end{equation}
With $aa^\dagger=N+1$ and $a^\dagger a=N$,
\begin{equation}
\boxed{H_{jk}=\hbar\omega(k+\tfrac12)\delta_{jk}}.
\end{equation}
\qnote{Truncating $X$ and $P$ before squaring
can corrupt their topmost diagonal level.}

\problem{3.40}{Classical-frequency motion for any oscillator state}
Let
$\ket{\Psi(t)}=\sum_{n=0}^\infty
c_n\ee^{-\ii\omega(n+1/2)t}\ket n$
with $\sum|c_n|^2=1$.
The oscillator ladder matrix elements give
\begin{equation}
\avg{a}(t)=\ee^{-\ii\omega t}
\sum_{n=0}^\infty\sqrt{n+1}\,c_n^*c_{n+1}.
\end{equation}
Set $S=\sum\sqrt{n+1}\,c_n^*c_{n+1}$
and write
$S=|S|\ee^{\ii\phi}$.
Since
$x=\sqrt{\hbar/(2m\omega)}(a+a^\dagger)$,
\begin{equation}
\boxed{\avg x(t)=C\cos(\omega t-\phi)},
\quad
\boxed{C=\sqrt{\frac{2\hbar}{m\omega}}|S|}.
\end{equation}
Equivalently
$C\ee^{-\ii\phi}=
\sqrt{2\hbar/(m\omega)}
\sum\sqrt{n+1}\,c_{n+1}^*c_n$.
For Problem 2.40, the real coefficients are
$c_0=3/5$, $c_1=-2\sqrt2/5$,
$c_2=2\sqrt2/5$, with all others zero.
Then
\begin{equation}
S=c_0c_1+\sqrt2c_1c_2
=-\frac{14\sqrt2}{25},
\end{equation}
and
\begin{equation}
\boxed{C=\frac{28}{25}
\sqrt{\frac{\hbar}{m\omega}},\quad \phi=\pi}.
\end{equation}
Thus every oscillator state with finite
$\avg x$ follows the \emph{classical
frequency} $\omega$, even if its
probability density is highly nonclassical.
